% ===================== Chapter 17 =====================
\chapter{The Standard Library, Part 2}
\label{ch:stdlib-2}

The first standard-library chapter covered computation: numbers, text, time. This
one covers \emph{structure and the outside world} the containers that organise
your data, and the modules that read files, parse data formats, and match
patterns. With these you can write programs that do real work on real data.

\section{Collections beyond the vector}

You met \code{Vector<T>} in Chapter~7. The \code{collections} module which is
available without an explicit import for the core containers offers several more
generic data structures, each suited to a particular access pattern.

\subsection{HashMap: key-value lookup}

A \code{HashMap<K, V>} associates \textbf{keys} of type \code{K} with
\textbf{values} of type \code{V}, and looks a value up by its key almost
instantly. It is the right tool whenever you think ``given \emph{this}, find
\emph{that}'' a username to a score, a word to its count, a product code to a
price.

\begin{salam}
func main:
    mut prices HashMap<str, i32> = HashMap {}
    prices.insert("apple", 3)
    prices.insert("banana", 5)
    println "size:", prices.size()
    println "apple costs:", prices.get("apple")
    println "have banana?", prices.has("banana")
    println "have cherry?", prices.has("cherry")
    prices.remove("apple")
    println "size now:", prices.size()
    prices.free()
end
\end{salam}

\begin{progout}
size: 2
apple costs: 3
have banana? true
have cherry? false
size now: 1
\end{progout}

The key methods are \code{insert(k, v)} to add or update an entry, \code{get(k)}
to retrieve a value, \code{has(k)} to test for a key, \code{remove(k)} to delete
one, \code{size()} for the count, and \code{free()} when you are done.

\begin{warn}
Call \code{get} only for a key you know is present check first with \code{has}
if you are not sure. As with vectors, a \code{HashMap} manages its own memory, so
remember \code{free()} when you have finished with it.
\end{warn}

\subsection{Stack and Queue}

When you care about the \emph{order} in which items come back out, two classic
containers help. A \code{Stack<T>} is last-in-first-out the most recently pushed
item pops first, like a stack of plates. A \code{Queue<T>} is first-in-first-out ---
items leave in the order they arrived, like a line at a counter.

\begin{salam}
func main:
    mut s Stack<i32> = Stack {}
    s.push(1)
    s.push(2)
    println "top:", s.peek()        // 2
    println "pop:", s.pop()         // 2 (last in, first out)
    s.free()

    mut q Queue<str> = Queue {}
    q.enqueue("first")
    q.enqueue("second")
    println q.dequeue()             // "first" (first in, first out)
    println q.dequeue()             // "second"
    q.free()
end
\end{salam}

\begin{progout}
top: 2
pop: 2
first
second
\end{progout}

A \code{Stack} uses \code{push}/\code{pop}/\code{peek}; a \code{Queue} uses
\code{enqueue}/\code{dequeue}/\code{peek}. Both report \code{size()} and
\code{is\_empty()}, and both should be \code{free}d. Reach for a stack when you
need to undo or backtrack (the most recent thing first), and a queue when you must
process things fairly in arrival order.

\section{Reading and writing files}

The \code{io} module, which you met for console output, also reads and writes
files. The simplest operations handle a whole file at once:

\begin{salam}
import "io"

func main:
    io.WriteFile("notes.txt", "line one\nline two\n")
    content auto = io.ReadFile("notes.txt")
    print content
end
\end{salam}

\begin{progout}
line one
line two
\end{progout}

\code{io.WriteFile(path, text)} writes (creating or replacing the file), and
\code{io.ReadFile(path)} returns the whole file as a \code{str}. For
line-by-line processing, \code{io.Lines(path)} returns a \code{Vector<str>}, one
entry per line, which you can walk with the techniques from Chapter~7.

\section{The operating system: \texttt{os}}

The \code{os} module is your program's window onto the system it runs on: files
and directories, environment variables, command-line arguments, and the exit
code.

\begin{salam}
import "os"

func main:
    println "exists?", os.Exists("notes.txt")
    println "current dir:", os.Cwd()
    home auto = os.Env("HOME")
    println "home:", home
end
\end{salam}

Common \code{os} operations include \code{os.Exists(path)} to test for a file,
\code{os.Remove(path)} to delete one, \code{os.Mkdir(path)} to create a directory,
\code{os.ListDir(path)} to list a directory's contents as a \code{Vector<str>},
\code{os.Args()} for the command-line arguments, \code{os.Env(name)} for an
environment variable, and \code{os.Exit(code)} to end the program with a status
code.

\begin{salam}
import (
    "os"
    "io"
)

func main:
    // back up a file, then clean up after ourselves
    io.WriteFile("data.txt", "important\n")
    defer os.Remove("data.txt")
    io.WriteFile("data.bak", io.ReadFile("data.txt"))
    defer os.Remove("data.bak")
    println "backup made; files exist?", os.Exists("data.bak")
end
\end{salam}

Notice the \code{defer}s from Chapter~14 doing exactly what they were made for:
each temporary file is scheduled for removal the instant it is created, so the
program tidies up no matter how it ends.

\section{Working with JSON}

JSON is the lingua franca of data exchange, and the \code{json} module reads it
without ceremony. You can validate a document, pull out fields by key, and ask
whether a key is present:

\begin{salam}
import "json"

func main:
    doc str = "{ \"user\": \"ada\", \"age\": 36 }"
    println "valid?", json.Valid(doc)
    println "user:", json.Get(doc, "user")
    println "age:", json.GetInt(doc, "age")
    println "has email?", json.Has(doc, "email")
end
\end{salam}

\begin{progout}
valid? true
user: ada
age: 36
has email? false
\end{progout}

\code{json.Valid} checks that the text is well-formed JSON; \code{json.Get} returns
a field as a string, while \code{json.GetInt} and \code{json.GetFloat} return it as
a number; \code{json.Has} tests for a key. There is also \code{json.Escape} to make
a string safe to embed in JSON, and \code{json.Minify} to strip whitespace. This is
enough to consume the JSON that web services and config files speak.

\section{Pattern matching with regular expressions}

A \textbf{regular expression} (regex) is a compact pattern that describes a set of
strings ``one or more digits,'' ``a word followed by a space.'' The
\code{regex} module matches text against such patterns. For one-off use, the
\code{...Str} functions take the pattern as a plain string:

\begin{salam}
import "regex"

func main:
    println "has digits?", regex.MatchStr("[0-9]+", "abc123")
    println "digits start at:", regex.FindStr("[0-9]+", "abc123def")
    println regex.ReplaceStr("[0-9]+", "a1b22c333", "#")
end
\end{salam}

\begin{progout}
has digits? true
digits start at: 3
a#b22c333
\end{progout}

\code{regex.MatchStr(pattern, text)} reports whether the pattern occurs anywhere in
the text; \code{regex.FindStr} returns the index where the first match begins (or a
negative value if there is none); and \code{regex.ReplaceStr} substitutes a match
with replacement text.

\begin{deep}[In Depth: compile once, match many]
If you will apply the same pattern many times say, validating every line of a
large file compiling it once is more efficient than re-parsing the pattern on
each call. \code{regex.Compile(pattern)} returns a reusable \code{Regex} value, and
\code{regex.Match(re, text)}, \code{regex.Find(re, text)}, and
\code{regex.Replace(re, text, with)} operate on it. The \code{...Str} convenience
functions are exactly these with an implicit compile each time; prefer the
compiled form in a hot loop.
\end{deep}

\section{Other data-format modules}

Rounding out the picture, the standard library also includes a \code{csv} module
for reading and writing comma-separated values (the format spreadsheets export),
an \code{encoding} module for common encodings, and the networking modules
\code{net}, \code{tcp}, and \code{http} for talking to other machines. The shape
is always the same as what you have seen here: import the module, call its
qualified functions, and free anything that manages its own memory. Once you are
comfortable with the modules in these two chapters, the rest follow the same
patterns.

\section{A worked example: word frequencies}

Let us combine a \code{HashMap}, the \code{str} module, and a vector into a small
but genuinely useful program: counting how often each word appears in a sentence.

\begin{salam}
import "str"

func main:
    text str = "the cat sat on the mat the cat ran"
    mut counts HashMap<str, i32> = HashMap {}
    mut words Vector<str> = str.Split(text, " ")
    mut i i32 = 0
    while i < words.len():
        w auto = words.get(i)[0]
        if counts.has(w):
            counts.insert(w, counts.get(w) + 1)
        else:
            counts.insert(w, 1)
        end
        i += 1
    end
    println "distinct words:", counts.size()
    println "the appears", counts.get("the"), "times"
    println "cat appears", counts.get("cat"), "times"
    words.free()
    counts.free()
end
\end{salam}

\begin{progout}
distinct words: 6
the appears 3 times
cat appears 2 times
\end{progout}

Split the sentence into words, then for each word either start its count at 1 or
add one to its existing count, using \code{has} to tell the two cases apart. The
\code{HashMap} does the heavy lifting instant lookup by word and the result
is a tally of the text. This ``count occurrences'' pattern appears everywhere,
from analytics to spell-checkers.

\section{Chapter summary}

\begin{itemize}[leftmargin=1.4em]
  \item \code{HashMap<K, V>} maps keys to values: \code{insert}, \code{get},
        \code{has}, \code{remove}, \code{size}, \code{free}.
  \item \code{Stack<T>} (push/pop/peek) is last-in, first-out;
        \code{Queue<T>} (enqueue/dequeue) is first-in, first-out; both
        provide \code{free}.
  \item \code{io.WriteFile}, \code{io.ReadFile}, and \code{io.Lines} handle whole
        files; \code{os} provides \code{Exists}, \code{Remove}, \code{Mkdir},
        \code{ListDir}, \code{Args}, \code{Env}, \code{Exit}.
  \item \code{json} reads JSON: \code{Valid}, \code{Get}, \code{GetInt},
        \code{Has}, \code{Escape}, \code{Minify}.
  \item \code{regex} matches patterns: \code{MatchStr}, \code{FindStr},
        \code{ReplaceStr}, or the compiled \code{Compile}/\code{Match}/\code{Find}.
  \item \code{csv}, \code{encoding}, and the networking modules round out the
        library, all following the same import-and-qualify pattern.
\end{itemize}

\section{Exercises}

\exercise Build a \code{HashMap<str, i32>} of three people's ages, then look up and
print one age and test for a name that is not present.

\exercise Use a \code{Stack} to reverse the order of the numbers 1 to 5: push them
all, then pop them into a line of output.

\exercise Write a program that writes three lines to a file with
\code{io.WriteFile}, reads them back with \code{io.Lines}, and prints how many
lines there were.

\exercise Given the JSON \code{"\{ \textbackslash"title\textbackslash":
\textbackslash"Salam\textbackslash", \textbackslash"pages\textbackslash": 400
\}"}, print the title and the page count.

\exercise Use \code{regex.MatchStr} to test whether a string looks like it contains
a year (four digits in a row).

\exercise Extend the word-frequency example to also print the word ``the'' count
only if it appears more than twice.
